\begin{activity}{Building Geological Intuition with Numerical Experiments}

\section*{Purpose}
To use numerical experiments to explore how thermal structure responds to
geological complexity.

{\color{red}
\section*{Learning Outcomes}

By the end of this activity, you should be able to:
\begin{itemize}
  \item Predict qualitatively how depth-dependent conductivity, layered heat production, time-varying basal heat flux, or a transient intrusion would change a model's temperature field.
  \item Distinguish which geological complexities can still be handled analytically and which require a numerical approach.
  \item Explain how a numerical model isolates physical effects that cannot be separated observationally.
\end{itemize}
}

Consider modifying a numerical heat flow model to include:
\begin{itemize}
  \item depth-dependent thermal conductivity,
  \item layered heat production,
  \item time-varying basal heat flux,
  \item transient magma intrusions.
\end{itemize}

\section*{Task 1: Qualitative predictions}

For each modification listed above, predict how the temperature field would respond compared to a uniform, steady-state reference model.
\answerbox[8cm]

\section*{Task 2: Analytic vs numerical}

Identify which of the modifications could still be handled with an analytic solution, and which require a fully numerical approach. Justify your choices.
\answerbox[5cm]

\section*{Task 3: Isolating effects}

Explain how a numerical model allows you to separate physical effects that are difficult to disentangle from field observations.

Give one geological example where this ability is particularly valuable.
\answerbox[5cm]

{\color{red}
\section*{Key Ideas}

Numerical modelling enables controlled experiments on Earth systems that cannot be performed in nature.

\textbf{Numerical models are controlled experiments.} Unlike field observations, numerical experiments allow one variable to be changed at a time while holding all others fixed. This is the only practical way to isolate the thermal effect of, for example, a single high-conductivity layer within a complex lithosphere.

\textbf{Analytic solutions fail when properties vary in space or time.} The steady-state heat equation has a closed-form solution for uniform conductivity and constant heat production. Depth-dependent $k$, layered $A$, or a time-varying basal flux break these assumptions and generally require numerical integration.

\textbf{Qualitative intuition guides model design.} Before running a model, predicting the sign and magnitude of an effect (will a high-conductivity layer raise or lower temperatures above it?) builds the physical intuition needed to detect numerical errors and interpret model output.
}

\end{activity}
